\chapter{An exact computation workbook}
\label{app:computation-workbook}

This appendix gives calculations that a reader can reproduce without trusting a status flag. The objects are the normalized coefficients of $\Delta$, the auxiliary shell polynomials, and one rational transfer certificate at $31$. Each calculation ends at its actual mathematical boundary. Agreement of finite outputs does not supply the missing all-prime nonvanishing argument.

The complete integer inputs and outputs are in the companion evidence archive. Paths below are relative to its extracted project root. A JSON integer remains an integer even when it has many digits. An encoded rational pair $[a,b]$ means $a/b$, with a nonzero denominator. Converting these objects to floating point before checking an equality changes the contract.

\section{Three engines and one coefficient shift}
\label{sec:workbook-engines}

Write
\[
 F(q)=\frac{\Delta(q)}q=\prod_{k\ge1}(1-q^k)^{24}
      =\sum_{m\ge0}f_mq^m.
\]
Then $f_0=1$ and $\tau(n)=f_{n-1}$. This shift is essential. A program that stores $\tau(0)=0$ as a sentinel must not confuse that entry with the constant coefficient $f_0$.

For direct product expansion, terms of degree at most $m$ depend only on factors with $k\le m$. The exact binomial expansion
\[
 (1-q^k)^{24}=\sum_{j=0}^{24}(-1)^j\binom{24}{j}q^{kj}
\]
therefore supplies a finite algorithm for each desired coefficient. Three short calculations are
\begin{align*}
 f_1&=-24,\\
 f_2&=\binom{24}{2}-24=276-24=252,\\
 f_3&=-\binom{24}{3}+(-24)(-24)-24\\
    &=-2024+576-24=-1472.
\end{align*}
The last line comes respectively from the $q^3$ term in the first factor, the product of its $q$ term with the second factor's $q^2$ term, and the third factor's $q^3$ term. Thus it computes $\tau(4)$, not $\tau(3)$.

The logarithmic derivative provides a different exact engine. With $\sigma_1(j)=\sum_{d\mid j}d$, comparison of coefficients in
\[
 q\frac{F'(q)}{F(q)}=-24\sum_{j\ge1}\sigma_1(j)q^j
\]
gives
\begin{equation}
 m f_m=-24\sum_{j=1}^{m}\sigma_1(j)f_{m-j}
 \qquad(m\ge1).
 \label{eq:workbook-divisor-recurrence}
\end{equation}
For $m=3$, the divisor sums are $1,3,4$. Hence
\[
 f_3=-8\bigl(252+3(-24)+4\bigr)=-8\cdot184=-1472.
\]
An implementation should check that the numerator in \eqref{eq:workbook-divisor-recurrence} is divisible by $m$, rather than rounding the quotient.

Jacobi gives a third engine. Put $T_j=j(j+1)/2$ and $a_j=(-1)^j(2j+1)$. The series
\[
 A(q)=\prod_{k\ge1}(1-q^k)^3=1+\sum_{j\ge1}a_jq^{T_j}
\]
satisfies $F=A^8$. The identity $AqF'=8FqA'$ yields
\begin{equation}
 m f_m=\sum_{\substack{j\ge1\\T_j\le m}}
          a_j(9T_j-m)f_{m-T_j}.
 \label{eq:workbook-triangular-recurrence}
\end{equation}
At $m=3$, only $T_1=1$, $T_2=3$ contribute. Its numerator is
\[
 (-3)(9-3)(252)+5(27-3)(1)=-4536+120=-4416.
\]
Division by $3$ again gives $-1472$. The direct product, divisor recurrence and triangular recurrence reach the same integer through different algebraic operations.

\Needspace{11\baselineskip}
\begin{center}
\begin{tabular}{r r r r}
\toprule
$m$ & $f_m=\tau(m+1)$ & $m$ & $f_m=\tau(m+1)$\\
\midrule
0 & $1$ & 5 & $-6048$\\
1 & $-24$ & 6 & $-16744$\\
2 & $252$ & 7 & $84480$\\
3 & $-1472$ & 8 & $-113643$\\
4 & $4830$ & 9 & $-115920$\\
\bottomrule
\end{tabular}
\end{center}

The current independent comparison covers every $\tau(n)$ for $1\le n\le512$. The historical final-arithmetic run's direct-product comparison stopped at $128$, while the packet-tree run compared through $512$. These are different documented ranges. The new three-engine table supplies its own bound and does not infer one from an older summary.

\section{Reconstructing the auxiliary polynomial}
\label{sec:workbook-polynomial}

For a packet $v$, set $c_i=(v_i-i)/5$ and form
\[
 P_v(T)=\prod_{0\le i<j\le4}\bigl((i-j)+T(c_i-c_j)\bigr).
\]
Its physical value is $P_v(5)=288W(v)$. Summing the full shell gives
\[
 P_n(T)=\sum_{v\in\Vn}P_v(T)=\sum_{j=0}^{10}B_jT^j,
 \qquad P_n(5)=288\tau(n).
\]
The polynomial has a different variable from the prime index. A root in $T$ or in a rescaled variable does not represent a nearby prime coefficient.

For a nonzero $P_n$, use the convention $v_5(0)=+\infty$ and define
\[
 \mu_n=\min_j\bigl(j+v_5(B_j)\bigr),\qquad
 G_n(U)=5^{-\mu_n}P_n(5U)=\sum_{j=0}^{10}C_jU^j.
\]
Every $C_j=B_j5^{j-\mu_n}$ is an integer. The power of $5$ should be removed only after verifying divisibility for each coefficient. Since $288$ is a $5$-adic unit,
\begin{equation}
 v_5(\tau(n))=\mu_n+v_5(G_n(1)).
 \label{eq:workbook-normalized-valuation}
\end{equation}
All coefficients in the next two tables are exact integers from the freshly reproduced shell data. They allow a reader to check the normalization and every displayed residue without re-enumerating a large shell.

\section{The simple branch at 499}
\label{sec:workbook-499}

At $n=499$ the shell has $500$ packets. Of their weights, $262$ are positive and $238$ negative. Their sum is
\[
 \tau(499)=-108877719272500.
\]
The minimum layer valuation is $\mu_{499}=3$.

\Needspace{17\baselineskip}
\begin{center}
\begingroup\small
\setlength{\tabcolsep}{5pt}
\begin{tabular}{r r r r}
\toprule
$j$ & $B_j$ in $P_{499}$ & $C_j$ in $G_{499}$ & $j+v_5(B_j)$\\
\midrule
0 & $144000$ & $1152$ & 3\\
1 & $-360000$ & $-14400$ & 5\\
2 & $1063040$ & $212608$ & 3\\
3 & $28418240$ & $28418240$ & 4\\
4 & $-148872808$ & $-744364040$ & 4\\
5 & $209474504$ & $5236862600$ & 5\\
6 & $2422451240$ & $302806405000$ & 7\\
7 & $16445000$ & $10278125000$ & 11\\
8 & $8107402848$ & $25335633900000$ & 8\\
9 & $-11893316304$ & $-185833067250000$ & 9\\
10 & $-1160632800$ & $-90674437500000$ & 12\\
\bottomrule
\end{tabular}
\endgroup
\end{center}

Summing the second coefficient column and using the normalization checks gives
\begin{align*}
 P_{499}(5)&=-31356783150480000=288\tau(499),\\
 G_{499}(1)&=-250854265203840.
\end{align*}
Reducing the normalized coefficients gives
\begin{align*}
 G_{499}(U)&\equiv2+3U^2\pmod5,\\
 G_{499}(U)&\equiv2+8U^2+15U^3+10U^4\pmod{25}.
\end{align*}
The first polynomial has the simple roots $1$ and $4$ modulo $5$. At $U=1$, the exact sum is $10$ modulo $25$ and its derivative is $1$ modulo $5$. Consequently
\[
 G_{499}(1+5h)\equiv10+5h\pmod{25}.
\]
The five possibilities are $10,15,20,0,5$ as $h$ runs from $0$ to $4$. The root in the branch $U\equiv1\pmod5$ therefore lifts to $U=16\pmod{25}$.

This lifted root is not the physical target. The target remains $U=1$, where $v_5(G_{499}(1))=1$. Formula \eqref{eq:workbook-normalized-valuation} gives $v_5(\tau(499))=3+1=4$.

The following exact root history records both simple branches. Each entry lies between $0$ and $5^k-1$ and evaluates to zero modulo $5^k$. Consecutive entries agree modulo the preceding power. These conditions are stronger than displaying an approximate root.

\Needspace{18\baselineskip}
\begin{center}
\begin{tabular}{r r r r}
\toprule
$k$ & $5^k$ & Root reducing to $1$ & Root reducing to $4$\\
\midrule
1 & 5 & 1 & 4\\
2 & 25 & 16 & 4\\
3 & 125 & 41 & 54\\
4 & 625 & 541 & 554\\
5 & 3125 & 3041 & 2429\\
6 & 15625 & 9291 & 14929\\
7 & 78125 & 9291 & 77429\\
8 & 390625 & 165541 & 233679\\
9 & 1953125 & 1337416 & 1014929\\
10 & 9765625 & 7196791 & 1014929\\
11 & 48828125 & 7196791 & 20546179\\
12 & 244140625 & 153681166 & 215858679\\
\bottomrule
\end{tabular}
\end{center}

The fixed degree $10$ does not impose a universal ceiling on the valuation of polynomial values at varying inputs. A simple $5$-adic root supplies arbitrarily accurate integer approximations. That observation concerns the auxiliary variable. It neither changes the value at $1$ nor supplies interpolation of the prime-indexed family.

\section{The double branch at 16249}
\label{sec:workbook-16249}

At $16249$ the shell contains $16250$ packets, with $8284$ positive and $7966$ negative weights. The exact coefficient is
\[
 \tau(16249)=133116215274958169206250.
\]
Here $\mu_{16249}=4$. The full coefficient table is as follows.

\Needspace{17\baselineskip}
\begin{center}
\begingroup\small
\setlength{\tabcolsep}{4pt}
\begin{tabular}{r r r r}
\toprule
$j$ & $B_j$ in $P_{16249}$ & $C_j$ in $G_{16249}$ & $j+v_5(B_j)$\\
\midrule
0 & $4680000$ & $7488$ & 4\\
1 & $-13320000$ & $-106560$ & 5\\
2 & $-199541600$ & $-7981664$ & 4\\
3 & $-26375464320$ & $-5275092864$ & 4\\
4 & $737398074300$ & $737398074300$ & 6\\
5 & $-39838758304992$ & $-199193791524960$ & 5\\
6 & $-746041848478320$ & $-18651046211958000$ & 7\\
7 & $-6990990388478280$ & $-873873798559785000$ & 8\\
8 & $21389931029028020$ & $13368706893142512500$ & 9\\
9 & $-64003552771660728$ & $-200011102411439775000$ & 9\\
10 & $3937759175521630080$ & $61527487117525470000000$ & 11\\
\bottomrule
\end{tabular}
\endgroup
\end{center}

The exact sums are
\begin{align*}
 P_{16249}(5)&=38337469999187952731400000,\\
 G_{16249}(1)&=61339951998700724370240.
\end{align*}
The first equals $288\tau(16249)$, and the second equals the first divided by $625$. Modulo $5$, the normalized polynomial is
\[
 3+U^2+U^3=(U-1)^2(U-2).
\]
Modulo $25$, the table gives
\[
 G_{16249}(U)\equiv13+15U+11U^2+11U^3+15U^5.
\]
Thus $G_{16249}(1)\equiv15\pmod{25}$ while $G'_{16249}(1)\equiv0\pmod5$. For every integer or $5$-adic integer $h$, polynomial Taylor expansion gives
\[
 G_{16249}(1+5h)\equiv G_{16249}(1)+5hG'_{16249}(1)
                    \equiv15\pmod{25}.
\]
Terms of order at least two have a factor $25$. This is a uniform obstruction throughout that residue branch, not merely the observation that five trial inputs fail. Every input congruent to $1$ modulo $5$ has normalized valuation exactly $1$.

The other residual root should not be discarded. At $U=2$, the normalized value is $5$ modulo $25$ and its derivative is $1$ modulo $5$. That simple branch lifts to $22$ modulo $25$. The conclusion is that the double branch at $1$ does not lift, rather than that the polynomial has no $5$-adic roots anywhere.

At the physical input $1$, the coefficient valuation is $4+1=5$. The two examples have different root geometry but both settle their actual target value by an exact modulo-$25$ calculation.

\section{A first failure needs a domain}
\label{sec:workbook-valuation-domain}

The tested auxiliary equality was
\[
 v_\ell(\tau(p))=v_\ell(p+1),
 \qquad p\ne\ell,\quad\ell\mid(p+1),\quad p\text{ prime}.
\]
The complete search through $150000$ contains $13848$ primes. The following counts refer to the indicated candidate sets, not to every prime without a residue condition.

\Needspace{11\baselineskip}
\begin{center}
\begingroup\small
\begin{tabular}{r r r r r r}
\toprule
$\ell$ & First $p$ & $v_\ell(\tau(p))$ & $v_\ell(p+1)$ & Candidates & Discrepancies\\
\midrule
2 & 81919 & 20 & 14 & 13847 & 1\\
3 & 8747 & 8 & 7 & 6941 & 5\\
5 & 499 & 4 & 3 & 3451 & 50\\
\bottomrule
\end{tabular}
\endgroup
\end{center}

The exact decisive values, including the discrepancy in the opposite $5$-adic direction, are
\begin{align*}
 \tau(499)&=-108877719272500,\\
 \tau(1249)&=27757672778617250,\\
 \tau(8747)&=348766170858826454196,\\
 \tau(81919)&=-904107824789465028682055680.
\end{align*}
For example, $1250=2\cdot5^4$, while the second integer has $5$-adic valuation $3$. Hence a proposed one-sided lift $v_5(\tau(p))\ge v_5(p+1)$ also fails at this actual prime.

These four coefficients were checked independently by the Eisenstein convolution. If $\sigma_r(n)=\sum_{d\mid n}d^r$, the coefficient of $(E_4E_8-E_6^2)/1728$ is obtained from
\begin{align*}
1728\tau(n)={}&240\sigma_3(n)+480\sigma_7(n)+1008\sigma_5(n)\\
 &+115200\sum_{k=1}^{n-1}\sigma_3(k)\sigma_7(n-k)\\
 &-254016\sum_{k=1}^{n-1}\sigma_5(k)\sigma_5(n-k).
\end{align*}
This computes the selected integers without the triangular recurrence. The numerator is divided exactly by $1728$. The replay checks this independent formula for $1\le n\le128$ and at the four displayed primes, giving $132$ checks.

The phrase ``first failure'' has two limits. It concerns the ordered search in its declared candidate set, and its exhaustive range ends at $150000$. A larger prime may supply another failure. The displayed nonzero integers refute an auxiliary valuation identity, not Lehmer's conjecture.

\section{Reading the rational transfer at 31}
\label{sec:workbook-transfer}

The degree-at-most-$22$ certificate has six coarse sectors and sixteen fine sectors. Both correction matrices have eighteen columns. Let
\[
 A\in\Z^{6\times16},\quad
 M_c\in\Z^{6\times18},\quad M_f\in\Z^{16\times18}.
\]
The fine-to-coarse parent map completely specifies $A$: its column $j$ has a $1$ in the parent row and zero elsewhere. The original fine masses are given below. Let $s_f$ denote their sixteen-entry column vector. Indices start at $0$.

\Needspace{15\baselineskip}
\begin{center}
\begin{tabular}{r r r @{\qquad} r r r}
\toprule
$j$ & Parent & Fine mass & $j$ & Parent & Fine mass\\
\midrule
0 & 0 & $2354352$ & 8 & 0 & $-11844448$\\
1 & 2 & $-19467448$ & 9 & 1 & $-7162848$\\
2 & 4 & $-832048$ & 10 & 0 & $-31156048$\\
3 & 4 & $15074752$ & 11 & 2 & $4356352$\\
4 & 4 & $1881152$ & 12 & 2 & $32821152$\\
5 & 5 & $4248552$ & 13 & 2 & $1262352$\\
6 & 5 & $7759752$ & 14 & 0 & $-7455448$\\
7 & 4 & $-31808448$ & 15 & 3 & $-12874848$\\
\bottomrule
\end{tabular}
\end{center}

For instance, row $0$ aggregates fine indices $0,8,10,14$:
\[
 2354352-11844448-31156048-7455448=-48101592.
\]
The six coarse masses are therefore
\[
 s_c=(-48101592,-7162848,18972408,
      -12874848,-15684592,12008304)^t.
\]
Their sum, and the sum of the fine masses, is $-52843168$. These original masses differ in order from the lexicographic orbit bars because the certificate orders the residue-sector keys.

The full machine certificate is
\begin{center}\small
\path{evidence/computation/replays/transport/p31_transport_certificate.json}.
\end{center}
It supplies an $18\times13$ rational matrix $N$ whose columns form a basis of $\ker M_c$, and the detail matrix $B=M_fN$ of size $16\times13$. The basic exact checks are
\[
 AM_f=M_c,\quad As_f=s_c,\quad M_cN=0,
 \quad B=M_fN,\quad AB=0.
\]
The last equation shows that a detail correction changes no coarse mass.

The rational row-reduction results are
\[
 \rank M_c=5,\qquad\rank M_f=15,\qquad\rank B=10.
\]
An independent check obtains the same lower bounds over both prime fields $\mathbb F_{65537}$ and $\mathbb F_{1000003}$. Every denominator used in a reduction is invertible there. A nonzero minor in either field proves a nonzero rational minor after clearing denominators.

The corresponding upper bounds have direct structural explanations. Each correction matrix has column sum zero, giving bounds $5$ and $15$. Every row of $A$ has a nonempty, disjoint support, so $\rank A=6$ and $\dim\ker A=10$. Since $AB=0$, the detail rank is at most $10$. The upper and lower bounds establish the displayed ranks exactly.

The coarse coefficient vector $a\in\Q^{18}$ makes
\[
 x=s_f+M_fa,\qquad Ax=(0,0,0,0,0,-52843168)^t.
\]
The certificate then chooses a coefficient vector $d\in\ker M_c$. Its final fine vector $y=x+M_fd$ is zero except at index $6$, where it equals $-52843168$.

One parent gives a concrete view of this second stage. Parent $5$ has fine children $5,6$. Just after the coarse stage their values are
\begin{align*}
 x_5&=-\frac{161598180629860124657534386744}
                 {23300823161076013875625},\\
 x_6&=-\frac{1069691132209170737342448593256}
                 {23300823161076013875625}.
\end{align*}
Their exact sum is $-52843168$. The detail stage adds $-x_5$ to child $5$ and $x_5$ to child $6$, giving $(0,-52843168)$ in that parent. All other parents retain their zero totals. This is a rational conservation calculation, not a numerical near-zero test.

To verify the whole certificate, parse the rational pairs, check the matrix dimensions, verify all five identities, and compare both stored stages entry by entry. Then check $M_cd=0$ and the single surviving final entry. The script \path{code/check_transport_certificate.py} performs the independent equality and prime-field checks.

This example does not assume a uniform rank theorem. Nor does full rank prove a uniformly nonzero survivor. At a shell whose original total were zero, complete zero-total transfer could still leave every final mass zero. At $31$ the separately computed exact total is nonzero, so the displayed survivor is authenticated for this one shell.

\section{Reproducing a run and interpreting its record}
\label{sec:workbook-manifests}

From the extracted source-set root, a new independent coefficient run is requested with
\begin{lstlisting}[basicstyle=\ttfamily\small,breaklines=true,frame=single,rulecolor=\color{Soft}]
python3 code/run_replays.py tau_independent \
  --output-root evidence/computation/reader_coefficients
\end{lstlisting}
The destination must be fresh. The launcher refuses to replace an existing run directory. Omitting the suite name requests the complete declared sequence of suites. For that complete reproduction, use a fresh destination such as \path{evidence/computation/reproduced}. The large modular comparison consumes the newly generated dyadic result in that same destination, so these dependent suites belong in their declared order.

A version $2$ manifest records the actual argument vector, exact arithmetic domain, semantic bounds, software versions and executed limits. Its \texttt{input\_artifacts} contain input checksums before and after execution. Its \texttt{declared\_results} name the scientific files the child must create. Output checksums cover those files and the execution logs.

Three questions remain separate. Did execution finish within its limits with unchanged inputs and all required outputs? Does the code implement the intended finite assertion? Does a proof show that the finite assertion decides the universal claim? An exit code of zero answers only part of the first question.

The count $704$ means $22$ selected modulus examples times $32$ exponents, not $704$ distinct primes. Of $24$ selected moduli, $11$ and $19$ had no observed matching prime through $10000$. That finite absence does not contradict positive density. The $24576$ synthetic cases mean $256$ moduli times three trace parameters times $32$ exponents. Their parameter $p=5M-1$ need not be prime, and their trace is not asserted to equal $\tau(p)$.

A checksum pins bytes. Run provenance records execution. A timeout records an execution limit. A reproduced obstruction can refute its stated auxiliary implication while leaving Lehmer unchanged. Attach each exact calculation to a claim of the correct scope.
